\section*{अध्याय \ref{ch:g2:shapes} --- आकार और ठोस}

\begin{solution}{exo:g2:shapes:1}
$5$ भुजाएँ, $5$ कोने: पंचभुज। $3$ कोने (और $3$ भुजाएँ):
त्रिभुज। गोल: वृत्त।
\end{solution}

\begin{solution}{exo:g2:shapes:2}
\omterm{def:g2:shapes:sort}{बहुभुज}: वर्ग, त्रिभुज, षट्भुज (सीधी भुजाएँ)। \omterm{def:g2:shapes:sort}{बहुभुज} नहीं:
वृत्त (वक्र सीमा)।
\end{solution}

\begin{solution}{exo:g2:shapes:3}
समकोण, जैसे: बोर्ड, किताब या खिड़की के कोने; समकोण नहीं, जैसे:
डिब्बे के खुले ढक्कन का कोण या $2$ बजे घड़ी की सुइयों के बीच
का कोण।
\end{solution}

\begin{solution}{exo:g2:shapes:4}
वर्ग: $4$ समकोण। आयत: $4$। त्रिभुज में ज़्यादा से ज़्यादा
$1$ समकोण --- दो होते तो बाक़ी दोनों भुजाएँ कभी मिलतीं ही
नहीं।
\end{solution}

\begin{solution}{exo:g2:shapes:5}
कोनों के बीच के खाने गिनकर बनी नक़लें; नमूने से मिलाने पर काम
की जाँच हो जाती है।
\end{solution}

\begin{solution}{exo:g2:shapes:6}
पासा: \omterm{def:g2:shapes:solids}{घन}। सूप का टिन: \omterm{def:g2:shapes:solids}{बेलन}। फुटबॉल: गेंद। जूतों का डिब्बा: डिब्बा।
\end{solution}

\begin{solution}{exo:g2:shapes:7}
\omterm{def:g2:shapes:solids}{घन} के $6$ \omterm{def:g2:shapes:solids}{फलक} होते हैं, हर एक वर्ग। डिब्बे के भी $6$ \omterm{def:g2:shapes:solids}{फलक} होते
हैं (आयत)।
\end{solution}

\begin{solution}{exo:g2:shapes:8}
गेंद लुढ़कती है (हर ओर से गोल) और \omterm{def:g2:shapes:solids}{बेलन} भी लुढ़कता है (अपनी वक्र
सतह पर)। \omterm{def:g2:shapes:solids}{घन} और डिब्बा नहीं लुढ़कते: उनके सब \omterm{def:g2:shapes:solids}{फलक} चपटे हैं और वे
उन्हीं पर रुक जाते हैं।
\end{solution}

\begin{solution}{exo:g2:shapes:9}
$6$ भुजाएँ: षट्भुज, जिसके $6$ कोने हैं।
\end{solution}

\begin{solution}{exo:g2:shapes:10}
घर की बाहरी रेखा (वर्ग वाली दीवार और ऊपर त्रिभुज वाली छत, जिसका
आधार वर्ग की ऊपरी भुजा पर बैठा है) में $5$ भुजाएँ और $5$ कोने
हैं --- यह पंचभुज है (नीचे की भुजा, दो दीवारें, छत की दो ढलानें;
दीवार का ऊपरी सिरा छत के नीचे छिप जाता है)।
\end{solution}
